\chapter*{Sommario}
\addcontentsline{toc}{chapter}{Sommario}

La guida autonoma rappresenta un ambito di ricerca particolarmente complesso, poiché richiede l'integrazione di percezione, pianificazione, decisione e controllo del veicolo in ambienti dinamici e non completamente prevedibili. Tra le situazioni più critiche che un veicolo autonomo deve affrontare vi è il superamento sicuro di un ostacolo presente lungo la corsia di marcia. Questa manovra non può essere gestita semplicemente come una deviazione dalla traiettoria, ma richiede una valutazione accurata del contesto circostante, della presenza di altri veicoli, dello spazio disponibile e della possibilità di effettuare un cambio corsia senza compromettere la sicurezza.

Questa tesi affronta il problema della progettazione, integrazione e valutazione di un sistema per la gestione del comportamento di un veicolo autonomo in ambiente simulato, con particolare attenzione al caso del rilevamento e superamento di ostacoli. Il lavoro è stato sviluppato a partire dal framework CARLA Leaderboard, utilizzato come infrastruttura per la simulazione e la valutazione dell'agente in scenari di guida realistici. Il progetto non si limita alla sola implementazione dell'agente, ma comprende anche l'adattamento del codice esistente, la configurazione dell'ambiente di simulazione, la gestione delle route e l'integrazione dei diversi moduli necessari all'esecuzione degli esperimenti.

Il sistema sviluppato si basa su un framework decisionale che analizza lo stato corrente del veicolo e dell'ambiente circostante per determinare la manovra più opportuna. In presenza di un ostacolo sulla traiettoria, il veicolo valuta se rallentare, arrestarsi oppure eseguire un cambio corsia. La possibilità di cambiare corsia viene determinata considerando la posizione dell'ostacolo, la distanza di sicurezza, la presenza di veicoli nelle corsie adiacenti e la disponibilità di spazio sufficiente per completare la manovra. In questo modo, il comportamento del veicolo non è basato soltanto sul semplice inseguimento di una route, ma tiene conto delle condizioni operative dello scenario.

Per quanto riguarda il controllo del veicolo, il lavoro considera l'utilizzo di controllori dedicati alla gestione della traiettoria e della velocità. I controllori PID sono impiegati per produrre comandi di guida reattivi e stabili, agendo su grandezze come sterzo, accelerazione e frenata. Il controllo MPC, invece, viene considerato come approccio per pianificare l'evoluzione della manovra su un orizzonte temporale, tenendo conto dei vincoli del veicolo e dell'ambiente. L'integrazione tra logica decisionale e controllo consente di collegare la scelta della manovra al comportamento effettivo del veicolo durante la simulazione.

La valutazione sperimentale è stata condotta in ambiente CARLA Leaderboard, sfruttando scenari simulati e route predefinite per osservare il comportamento dell'agente in presenza di situazioni complesse. Gli esperimenti hanno permesso di analizzare la capacità del sistema di riconoscere condizioni critiche, prendere decisioni coerenti con il contesto e mantenere un controllo stabile del veicolo durante la guida. I risultati ottenuti evidenziano l'importanza di un approccio integrato, in cui la gestione degli ostacoli, la valutazione della sicurezza del cambio corsia e il controllo del veicolo cooperano per migliorare l'affidabilità complessiva della guida autonoma.

Nel complesso, la tesi mostra come un sistema di guida autonoma possa essere sviluppato e valutato all'interno di un framework simulativo strutturato, evidenziando sia il ruolo dell'agente autonomo sia il lavoro di integrazione e adattamento del codice necessario per rendere possibile la sperimentazione. Il lavoro svolto costituisce una base per ulteriori sviluppi futuri, orientati al miglioramento della robustezza del sistema, all'estensione degli scenari considerati e alla valutazione del comportamento del veicolo in condizioni di traffico sempre più complesse.